\chapter{Esquema de Dados e Registros de Execução}\label{ap:esquema-dados}

Este apêndice documenta a organização de diretórios e os esquemas de dados adotados na base experimental \texttt{database/dataset\_obi\_senior/}, bem como o formato dos registros acumulativos em JSON e das planilhas consolidadas em CSV gerados durante a execução dos experimentos em \texttt{results/dataset\_obi\_senior/}.

\section{Estrutura da Base de Questões e Esquema \texttt{problem.json}}\label{ap:esquema-problem}

Na base \texttt{database/dataset\_obi\_senior/}, cada um dos 125 problemas do Nível Sênior da OBI (2016--2025) é armazenado em um diretório próprio nomeado com o título oficial da questão (ou acrescido do ano em caso de homonímia entre edições, como \texttt{Chuva\_2019} e \texttt{Chuva\_2022}). A estrutura interna de cada diretório organiza-se em quatro elementos padronizados:

\begin{itemize}
    \item \texttt{problem.json}: arquivo principal com o enunciado estruturado, especificações de entrada e saída, restrições, exemplos, figuras associadas e metadados da prova;
    \item \texttt{imgs/}: subdiretório contendo os arquivos de imagem (\texttt{.png} ou \texttt{.jpg}) extraídos do caderno oficial, presente nas 46 questões multimodais;
    \item \texttt{solutions/}: subdiretório com as soluções humanas de referência disponibilizadas pela organização da OBI (\texttt{.cpp}, \texttt{.c}, \texttt{.py}, \texttt{.java} ou \texttt{.js}), utilizadas na validação em modo oráculo;
    \item \texttt{test\_cases/inputs/} e \texttt{test\_cases/outputs/}: subdiretórios contendo os pares de arquivos de entrada (\texttt{.in}) e saída esperada (\texttt{.out}) dos 4.681 casos de teste oficiais.
\end{itemize}

A Tabela~\ref{tab:ap-esquema-problem} descreve os campos que compõem o arquivo \texttt{problem.json} validado pelo modelo de dados \texttt{Problem} (\texttt{src/models/problem.py}).

\begin{table}[!htbp]
\centering
\caption[Campos do esquema problem.json da base OBI Sênior]{Descrição dos campos que compõem o esquema \texttt{problem.json} de cada questão da base.}
\label{tab:ap-esquema-problem}
\small
\resizebox{\textwidth}{!}{%
\begin{tabular}{@{}llp{9.2cm}@{}}
\toprule
\textbf{Campo} & \textbf{Tipo} & \textbf{Descrição} \\
\midrule
\texttt{title}             & \texttt{string}         & Título oficial do problema no caderno de prova da OBI. \\
\texttt{statement}         & \texttt{string}         & Texto principal do enunciado, contendo marcadores \texttt{[nome\_img]} na posição original das figuras. \\
\texttt{input}             & \texttt{string}         & Especificação do formato dos dados fornecidos na entrada padrão (\texttt{stdin}). \\
\texttt{output}            & \texttt{string}         & Especificação do formato exigido na saída padrão (\texttt{stdout}). \\
\texttt{constraints}       & \texttt{string}         & Restrições matemáticas e intervalos numéricos das variáveis de entrada. \\
\texttt{examples}          & \texttt{list[object]}   & Lista de pares \texttt{\{"input": ..., "output": ...\}} com os exemplos públicos do enunciado. \\
\texttt{imgs}              & \texttt{list[string]}   & Lista com os nomes dos arquivos de imagem localizados no subdiretório \texttt{imgs/}. \\
\texttt{imgs\_description} & \texttt{object / list}  & Mapeamento de cada imagem para sua descrição textual estruturada, empregado na modalidade \modalidade{description}. \\
\texttt{rating}            & \texttt{list[int]}      & Pontuação das submissões ou subtarefas previstas no caderno oficial (ex.: \texttt{[100]}). \\
\texttt{year}              & \texttt{string}         & Ano da edição da OBI à qual a questão pertence (2016 a 2025). \\
\texttt{level}             & \texttt{string}         & Nível da modalidade Programação (\texttt{"senior"}). \\
\texttt{period}            & \texttt{string}         & Etapa seletiva da competição (\textit{Fase 1 / Local}, \textit{Fase 2 / Estadual} ou \textit{Fase 3 / Nacional}). \\
\texttt{topics}            & \texttt{list[string]}   & Tópicos algorítmicos associados à questão segundo a ementa da OBI. \\
\texttt{time\_limit}       & \texttt{float}          & Tempo máximo de execução permitido por caso de teste, em segundos. \\
\texttt{memory\_limit}     & \texttt{int}            & Consumo máximo de memória RAM permitido durante a execução, em megabytes (MB). \\
\bottomrule
\end{tabular}%
}
\end{table}

A seguir, apresenta-se um exemplo real de arquivo \texttt{problem.json} referente à questão \textit{Acelerador de partículas} (OBI 2020, Fase Local):

\begin{small}
\begin{Verbatim}[breaklines, breakanywhere, fontsize=\footnotesize, frame=single, framesep=2mm, rulecolor=\color{black!40}]
{
    "title": "Acelerador de partículas",
    "statement": "A universidade está inaugurando um grande acelerador de partículas, com um emissor e três sensores, numerados 1, 2 e 3. [...] A figura mostra o caminho por onde as partículas trafegam, com uma graduação de 1 quilômetro. Por exemplo, do emissor até o acelerador são 3 quilômetros e a circunferência do acelerador tem 8 quilômetros.\n[1.png]\nNeste problema, será dada a distância total, em quilômetros, percorrida por uma certa partícula trafegando do emissor até algum sensor e seu programa deve determinar qual sensor foi atingido pela partícula.",
    "input": "A entrada consiste de apenas uma linha contendo um inteiro D, representando a distância total percorrida pela partícula.",
    "output": "Seu programa deve imprimir uma linha contendo um inteiro, representando o número do sensor que a partícula atingiu.",
    "constraints": "6 ≤ D ≤ 800008\nD sempre será a distância total percorrida entre o emissor e algum sensor.",
    "examples": [
        { "input": "23", "output": "2" },
        { "input": "6", "output": "1" },
        { "input": "9192", "output": "3" }
    ],
    "imgs": ["1.png"],
    "rating": [100],
    "year": "2020",
    "level": "senior",
    "period": "Fase Local",
    "topics": ["matemática", "aritmética modular", "simulação"],
    "time_limit": 5.0,
    "memory_limit": 256,
    "imgs_description": {
        "1.png": "[Descrição da imagem: Diagrama esquemático do acelerador de partículas. À esquerda, há um emissor retangular do qual parte um tubo reto horizontal em direção ao leste com comprimento de 3 km, marcado com traços a cada 1 km. No ponto de 3 km, o tubo conecta-se tangencialmente à parte inferior de um acelerador circular de 8 km de circunferência [...]]"
    }
}
\end{Verbatim}
\end{small}

\section{Estrutura dos Registros Acumulativos de Execução (\texttt{.json})}\label{ap:esquema-execucao}

Durante as baterias experimentais, o módulo \texttt{Orchestrator} armazena o histórico completo de cada submissão em arquivos JSON organizados hierarquicamente pelo caminho \texttt{results/dataset\_obi\_senior/<modelo>/<prompt>/<linguagem>/<modalidade>/<questao>.json}. Cada arquivo contém uma lista de objetos JSON em que cada elemento registra uma tentativa de resolução e seu respectivo julgamento local, conforme especificado na Tabela~\ref{tab:ap-esquema-execucao}.

\begin{table}[!htbp]
\centering
\caption[Campos do registro JSON acumulativo de execução]{Estrutura de campos dos arquivos JSON gerados por questão em cada bateria experimental.}
\label{tab:ap-esquema-execucao}
\small
\resizebox{\textwidth}{!}{%
\begin{tabular}{@{}llp{9.2cm}@{}}
\toprule
\textbf{Campo} & \textbf{Tipo} & \textbf{Descrição} \\
\midrule
\texttt{attempt}          & \texttt{int}    & Índice sequencial da tentativa registrada para a questão naquela configuração. \\
\texttt{timestamp}        & \texttt{string} & Data e horário UTC da conclusão da requisição no padrão ISO~8601. \\
\texttt{model}            & \texttt{string} & Identificador do modelo consultado na API (ex.: \texttt{gemini-3.8-flash}). \\
\texttt{prompt\_name}     & \texttt{string} & Nome do modelo de \textit{prompt} aplicado (\texttt{zero\_shot}, \texttt{few\_shot\_c} ou \texttt{chain\_of\_thought\_polya}). \\
\texttt{prompt}           & \texttt{string} & Texto integral do \textit{prompt} formatado enviado ao modelo na requisição. \\
\texttt{language}         & \texttt{string} & Linguagem-alvo da submissão (\texttt{python}, \texttt{c}, \texttt{cpp}, \texttt{java} ou \texttt{javascript}). \\
\texttt{modality}         & \texttt{string} & Modalidade de representação visual adotada (\texttt{text}, \texttt{img} ou \texttt{description}). \\
\texttt{include\_limits}  & \texttt{bool}   & Indica se os limites de tempo e memória foram injetados explicitamente no \textit{prompt}. \\
\texttt{time\_limit}      & \texttt{float}  & Limite de tempo de execução configurado no \texttt{JudgeService} (em segundos). \\
\texttt{memory\_limit}    & \texttt{int}    & Limite de memória configurado no \texttt{JudgeService} (em MB). \\
\texttt{tokens}           & \texttt{object} & Contagem de \textit{tokens} da chamada: \texttt{input\_tokens}, \texttt{output\_tokens} e \texttt{total\_tokens}. \\
\texttt{cost\_usd}        & \texttt{object} & Custo financeiro em dólares (USD): \texttt{input\_cost}, \texttt{output\_cost} e \texttt{total\_cost}. \\
\texttt{raw\_response}    & \texttt{string} & Resposta textual bruta retornada pela API do modelo. \\
\texttt{extracted\_code}  & \texttt{string} & Código-fonte limpo extraído do bloco Markdown para compilação e teste. \\
\texttt{judge\_result}    & \texttt{object} & Resultado da avaliação local contendo o veredito final (\texttt{judge\_predict}), o maior tempo de execução entre os testes (\texttt{execution\_time}), a contagem por veredito (\texttt{counts}) e o total de casos avaliados (\texttt{total\_test\_cases}). \\
\bottomrule
\end{tabular}%
}
\end{table}

O trecho a seguir exemplifica um registro JSON produzido para a questão \textit{Acelerador de partículas} em linguagem \linguagem{C}:

\begin{small}
\begin{Verbatim}[breaklines, breakanywhere, fontsize=\footnotesize, frame=single, framesep=2mm, rulecolor=\color{black!40}]
[
    {
        "attempt": 1,
        "timestamp": "2026-11-18T02:26:26.250617+00:00",
        "model": "deepseek/deepseek-v4.1-flash",
        "prompt_name": "chain_of_thought_polya",
        "prompt": "Você é um programador competitivo experiente em C. [...]",
        "language": "c",
        "modality": "text",
        "include_limits": false,
        "time_limit": 5.0,
        "memory_limit": 256,
        "tokens": {
            "input_tokens": 1143,
            "output_tokens": 3226,
            "total_tokens": 4369
        },
        "cost_usd": {
            "input_cost": 3e-06,
            "output_cost": 0.001936,
            "total_cost": 0.001939
        },
        "raw_response": "## 1. COMPREENSÃO DO PROBLEMA [...] ```c\n#include <stdio.h>\nint main(void) { ... }\n``` [...]",
        "extracted_code": "#include <stdio.h>\n\nint main(void) {\n    long long D;\n    if (scanf(\"%lld\", &D) != 1) return 0;\n    int r = (int)(D % 8);\n    if (r == 6) printf(\"1\\n\");\n    else if (r == 7) printf(\"2\\n\");\n    else printf(\"3\\n\");\n    return 0;\n}",
        "judge_result": {
            "judge_predict": "AC",
            "execution_time": 0.1003995,
            "counts": { "AC": 39, "WA": 0, "RE": 0, "TLE": 0, "MLE": 0, "CE": 0 },
            "total_test_cases": 39
        }
    }
]
\end{Verbatim}
\end{small}

\section{Formato das Tabelas CSV Consolidadas}\label{ap:esquema-csv}

Além dos registros detalhados em JSON, o ambiente persiste tabelas em formato CSV para sumarização direta dos vereditos e da telemetria em \texttt{results/dataset\_obi\_senior/}. Existem dois padrões de arquivos tabulares:

\begin{enumerate}
    \item \textbf{Tabelas de julgamento (\texttt{oracle.csv} e \texttt{judge/*.csv}):} consolidam o resultado do \texttt{JudgeService} por questão em cada bateria experimental. Cada linha corresponde a uma submissão avaliada e contém dez colunas:
    \begin{itemize}
        \item \texttt{question\_name}: identificador da questão na base;
        \item \texttt{execution\_time}: tempo máximo de execução (em segundos) registrado entre os casos de teste da questão;
        \item \texttt{judge\_predict}: veredito final atribuído à submissão (\veredito{AC}, \veredito{WA}, \veredito{RE}, \veredito{TLE}, \veredito{MLE} ou \veredito{CE});
        \item \texttt{AC}, \texttt{WA}, \texttt{RE}, \texttt{TLE}, \texttt{MLE}, \texttt{CE}: quantidade de casos de teste individuais classificados em cada categoria de veredito;
        \item \texttt{total\_test\_cases}: número total de casos de teste oficiais aferidos na questão.
    \end{itemize}
    \item \textbf{Tabela de comparação entre gerações (\texttt{llms\_compare.csv}):} reúne, para as 46 questões com figuras avaliadas na PP5 em \linguagem{Python} \textit{Zero-Shot}, a coluna \texttt{question\_name} seguida pelas triplas de colunas \texttt{veredito\_<modelo>}, \texttt{tokens\_<modelo>} (soma de \textit{tokens} da requisição) e \texttt{custo\_<modelo>} (custo em USD) para os modelos das gerações anterior e atual sob paridade de juiz.
\end{enumerate}
